\chapter
[\CO{[PKU]} Physikalisch-körperliche Untersuchung]
{\CC{[PKU]}\\Physikalisch-körperliche Untersuchung}
\label{chp:PKU}

\SectionUnderDevelopment

\minitoc

\section{Grundlagen}

Orientierung/Orientierungslinien:
ventral: vordere Medianlinie, Sternal-Parasternalinie, Medioklavikularlinie (MCL) rechts und
links, vordere/mittlere Axilarlinie
horizontal: Interkostalräume, Rippen (1. tastbare Rippe unter der Klavikel ist die 2. Rippe),
Rippenbogen
dorsal: hintere Axillarlinie, Scapularlinie (Ang. inf. Scapulae), Vertebral- Paravertebralinie,
hintere Medianlinie; Zählen von Wirbelkörpern und Dornfortsätzen bezogen auf den 7. HWK
und den 11. BWK

\subsection{Inspektion}

\subsection{Palpation}

\subsection{Perkussion}

Schallqualitäten, Klopfschall (KS):
• Schenkelschall: leise (kleine Amplitude), hell (höher frequent), kurz über nicht lufthaltigen
Organen; z.\,B. Leber und großen Ergüssen
• sonorer Klopfschall: über lufthaltigen Organen (Lungenschall) laut (große
• Amplitude), tief (niedere Frequenz), lang
• hypersonorer Klopfschall: bei Emphysem, Pneumothorax
• tympanitischer Klopfschall: (niederfrequent, eher rhythmische Schwingungen), lang,
• laut über gespannten, luftgefüllten Hohlorganen z.\,B. auch beim Gesunden über dem
Magenfundus, aber auch über dem Areal eines Pneumothorax / Spannungspneumothorax
(Schachtelton)

Orientierung/Orientierungslinien:
ventral: vordere Medianlinie, Sternal-Parasternalinie, Medioklavikularlinie (MCL) rechts und
links, vordere/mittlere Axilarlinie
horizontal: Interkostalräume, Rippen (1. tastbare Rippe unter der Klavikel ist die 2. Rippe),
Rippenbogen
dorsal: hintere Axillarlinie, Scapularlinie (Ang. inf. Scapulae), Vertebral- Paravertebralinie,
hintere Medianlinie; Zählen von Wirbelkörpern und Dornfortsätzen bezogen auf den 7. HWK
und den 11. BWK



\subsection{Auskultation}

Atemgeräusche:
- Vesikuläratmen: normales Atemgeräusch, Inspiration ist lauter als Exspiration
- Bronchialatmen: wie über Trachea oder oberen Lungenpartien mit großen Bronchien.
In- und Exspiration sind verschärft, z.\,B. über verdichtetem Lungengewebe
- trockene Rasselgeräusche(RG): Pfeifen und Giemen(expirat. Pfeifen)
Brummen
- feuchte RG: unterteilt in feinblasig (eher Knistern)und grobblasig (Blubbern),
in ohrnah und ohrfern,
in klingend und nicht klingend
Atemexkursion:
Feststellung nicht seitengleicher Beatmung: Auflegen der Hände und Beobachtung, ob sie
sich beim Atmen seitengleich verschieben
Stimmfremitus:
palpatorische Erfassung feiner Vibrationen mit den Fingerspitzen über den Lungenpartien
wenn der Patient mit tiefer Stimme 99 spricht; bei vergleichender Palpation verstärkte
Vibration über Infiltrationen, vermindert über Ergüssen und Pleuraschwarten
Perkussion:
Bronchophonie: Patient flüstert 66: bei Abhören mit Stethoskop ist das Geflüsterte
besser hörbar über verdichtetem Lungengewebe, Pneumonie, bei Infarkt, Atelektase,
vermindert über Pleuraerguss, -schwarte.


\section{Allgemeine Inspektion, Kopf und Halsuntersuchung}
Vorwissen 
Adipositas
Kachexie
Hautturgor
Exsikkose
Somnolenz
Sopor
pyknisch
athletisch
leptosom
Foetor ex ore
Exophtalmus
Enophtalmus
Myosis
Kayser-Fleischer Kornealring
Arcus senilis
Amaurosis
Bauchglatze
Suffusion
Quincke Ödem
Mydriasis
Maskengesicht
Akromegalie
Vollmondgesicht
Myxödem
Meningismus
Alopezie
Hirsutismus
Xanthelasmen
Leukoplakie
Struma nodosa
Exanthem
Ekzem
Petechien
Ikterus
Spider naevi
Angina tonsillaris
Anisokorie
Horner'sche Trias
Morbus Addison
Morbus Basedow

\subsection{Kurzer theoretischer Überblick}

Ein ausgeprägter Kropf (Struma) ist eventuell schon bei
zurückgeneigtem Kopf an der Assymetrie der Halses zu
erkennen. Die Schilddrüse (SD) ist bei Normalbefund in Ruhe
nicht tastbar. Bei der bimanuellen SD-Palpation steht der
Untersucher hinter dem sitzenden Patienten. Der Patient hält
den Kopf gerade und die Daumen des Untersuchers kommen im
Nacken zu liegen, die übrigen Finger palpieren suchend vom
Ringknorpel nach lateral. Dann legt der Untersucher die
Finger flach auf. Der Patient wird gebeten zu schlucken
(event. mit Hilfe eines Glas Wassers), die SD gleitet unter
den Fingern der Untersuchers vorbei. Pathologische Befunde
sind im Gegensatz zur regelrechten SD relativ leicht zu
tasten, z.B: Knoten, Vergrößerung, Assymetrie, fehlende
Schluckverschieblichkeit. Lymphknoten (LK) sind nur bei
Vergrößerung tastbar. Es handelt sich also um eine suchende
Palpation der entsprechenden Regionen. LK können reaktiv
entzündlich vergrößert sein, im Rahmen von malignen
Erkrankungen (Lymphom, Metastasen) oder systemischen
Erkrankungen (Sarkoidose etc). Palpation der axillären LK:
Der Patient lässt die Arme locker am Körper hängen. Der
Untersucher schiebt die Hände mit geschlossenen Fingern tief
in die Axilla und palpiert dann mit mäßigem Druck gegen die
Rippen nach caudal. Die Untersuchung wird weiter ventral und
dorsal wiederholt.

\subsection{Untersuchungsgang}



Untersuchungsgang



\Layout{2}{1.
Allgemeine
Inspektion}{%


Beurteilung von
Allgemeinzustand

wichtig: erster Eindruck des Patienten! Patient/In erscheint gesund – krank –
sehr krank? dyspnoisch, cyanotisch, blass, schweißig-kollaptisch, schockiert,
reduziert, sehr reduziert
• Tonus, Körperhaltung
• psychisch: normal, nervös, verwirrt
• Mimik: z.\,B. starr bei M. Parkinson, Tremor
• Sprache: normal, heiser, verwaschen
• Konstitutionstyp: athletisch, pyknisch, leptosom



Beurteilung von
Ernährungszustand

normal, reduziert, kachektisch, adipös

Größe/Gewicht: normal-, über-, untergewichtig; kachektisch
(BMI /Body Mass Index = kg/ Körpergröße in m2):
• Untergewicht: < 19
• Normalgewicht: (19) 20 – 25
• Übergewicht: 25 - 30
• Adipositas: 30 - 40
• Adipositas permagna: > 40
AZ und EZ, wenn nicht normal, sollen differenziert beschrieben werden


Beurteilung von
Bewusstseinslage

- Bewusstseinslage: „verlangsamt“, somnolent, soporös (schwer erweckbar),
bewusstlos, komatös

Beurteilung der Vitalparameter

normale Atemfrequenz: 12 - 20 /min
Tachypnoe: > 20/ min
Bradypnoe: < 12 /min
Atemnot – Patient/In ringt um Luft – dyspnoisch, orthopnoisch
Stridor: hörbares Geräusch bei behinderter Atmung durch Verengung der oberen
Luftwege,
Puls: normal 60-100/min, Tachykardie > 100 /min, Bradykardie < 60 /min
rhythmisch, arrhythmisch (Extrasystolen, Vorhofflimmern)
Blutdruck: Messung bei Erstkontakt notwendig!


der Kopfform, Gesichtsfarbe und
Mimik

Kopfform: symmetrisch, proportioniert
Beispiele für auffällige Krankheitsbilder: Akromegalie, Cushing Syndrom
(Vollmondgesicht), Tumore, Lähmungen
Gesichtsfarbe:
• Blässe: (Haut, Schleimhäute): Anämie, Kollaps, Nierenerkrankungen,
Myxödem
• Rotfärbung: Fieber, Cushing Syndrom, chron. Alkoholismus (Konjunktiven,
• Rhinophym), Hypertonus,
• Blaurotfärbung (Zyanose): periphere Zyanose, Polycythämie, Polyglobulie,
zentrale Zyanose (Vitium)
• Gelbfärbung: Ikterus (Skleren, Bauchhaut)
• Pigmentierung: Mb. Addison (Mundschleimhaut), Leberzirrhose,
Schwangerschaft, Medikamente, dermatologische Erkrankungen
}{}{}{}{}

\Layout{2}{2.
Inspektion}{%

\begin{description}
  \item[Augenmotorik]
Exophthalmus, Enophthalmus (Horner), Schielen, Lidödem, Lideinlagerungen
(Xanthelasmen)
Konjunktiven: blass, gerötet, blaurot, Blutungen, Verfärbungen
Skleren: Gefäßzeichnung, Gelbfärbung
Kornea: Kayser-Fleischer Kornealring (M. Wilson), Arcus senilis
\item[Pupillenform und -grösse]
\item[Lichtreaktion im
Seitenvergleich]
\item[Mund und
Rachen mit dem Spatel]
\item[Lippen]z.\,B. Bläschen, Einrisse, Karzinome, Schiefstellung, Ang. infectiosus
Mundschleimhaut: Bläschen, Soor, Leukoplakien, Blutungen
\item[Zunge]Beweglichkeit oder Atrophie (N. hypoglossus-Lähmung), Beläge,
Form- und Farbveränderung (atrophe Glossitis, chron.
Darmerkrankungen, AB Gabe,....), Tumore

\item[Gebisses]
\end{description}

}{}{}{}{}

\Layout{2}{3.
Palpation}{%

\begin{description}
  \item[Druckschmerzen](Tragusschmerz,
Trigeminusaustrittspunkte,
NNH)
\item[Schilddrüse](falls nötig beim Wassertrinken)

unauffällige Schilddrüse nicht sichtbar und tastbar;
Schilddrüse: Struma diffusa, Struma nodosa, Struma diffusa et nodosa

\item[LK-Status]des Halses
\end{description}
}{}{}{}{}

4.
Dokumentation




\section{Lungen- und Thoraxuntersuchung}

Vorwissen 

Dyspnoe
Tachypnoe
Bradypnoe
Kußmaulatmung
Cheyne-Stokes-Atmung
Biot'sche Atmung
Stridor
Blue bloater
Pink puffer
Emphysem
Atelektase
Pleuraerguß
Ellis-Damoiseau'sche Linie
Krönig'sche Schallfelder
Pneumothorax/Spannungpneu
COPD
Asthma bronchiale
Fassthorax
Lungenödem
Bronchiektasie
Lungenembolie
Silent Chest
Hyperventilation
Kyphose
Lordose
Skoliose
Vitalkapazität
Pleuraschwarte
Atypische Pneumonie

\subsection{Kurzer theoretischer Überblick}

Orientierung/Orientierungslinien:
ventral: vordere Medianlinie, Sternal-Parasternalinie, Medioklavikularlinie (MCL) rechts und
links, vordere/mittlere Axilarlinie
horizontal: Interkostalräume, Rippen (1. tastbare Rippe unter der Klavikel ist die 2. Rippe),
Rippenbogen
dorsal: hintere Axillarlinie, Scapularlinie (Ang. inf. Scapulae), Vertebral- Paravertebralinie,
hintere Medianlinie; Zählen von Wirbelkörpern und Dornfortsätzen bezogen auf den 7. HWK
und den 11. BWK

\Layout{2}{Schallqualitäten, Klopfschall (KS)}{%
\begin{description}
  \item[Schenkelschall]: leise (kleine Amplitude), hell
  (höher frequent), kurz über nicht lufthaltigen
  Organen; z.\,B. Leber und großen Ergüssen
  \item[sonorer Klopfschall]: über lufthaltigen Organen
  (Lungenschall) laut (große
  Amplitude), tief (niedere Frequenz), lang
  \item[hypersonorer Klopfschall]: bei Emphysem, Pneumothorax
  \item[tympanitischer Klopfschall]: (niederfrequent, eher
  rhythmische Schwingungen), lang, laut über gespannten,
  luftgefüllten Hohlorganen z.\,B. auch beim Gesunden über
  dem Magenfundus, aber auch über dem Areal eines
  Pneumothorax / Spannungspneumothorax
  (Schachtelton)
  \item[Dämpfung] des normalen sonoren Schalls (gedämpfter
  Klopfschall) bei Gewebsverdichtungen: 
  
  z.\,B. bei Pneumonie
  (pneumonische Infiltrationen werden nur erreicht, wenn
  sie nicht tiefer als etwa 5 cm unter der Oberfläche
  liegen und eine gewisse Ausdehnung erreichen) 
  
  z.\,B. bei
  Pleuraerguss und Pleuraschwarte (bei Erguss steigt die
  dorsale Obergrenze der Dämpfung von medial nach lateral hin
  an (Ellis-Damoiseau'sche Linie))
  
  z.\,B. bei Atelektase (bei verringertem Luftgehalt des
  Lungengewebes)
\end{description}

Je stärker geklopft wird, desto weniger sind kleinere
Veränderungen (dünner Erguss, kleinere Infiltrationen) zu
erfassen. Bei pathologischen Veränderungen der Lunge ist es
bei (seiten-) vergleichender Perkussion manchmal nicht
leicht festzustellen ob Unterschiede durch hypersonoren
Klopfschall der einen oder Dämpfung der anderen Seite
verursacht sind.
}{}{}{}{}

\Layout{2}{Atemgeräusche}{%
\begin{description}
  \item[Vesikuläratmen]: normales Atemgeräusch, Inspiration ist lauter als Exspiration
  \item[Bronchialatmen]: wie über Trachea oder oberen
  Lungenpartien mit großen Bronchien. In- und Exspiration
  sind verschärft, z.\,B. über verdichtetem Lungengewebe
  \item[trockene Rasselgeräusche(RG)]: Pfeifen und
  Giemen(expirat. Pfeifen), Brummen
  \item[feuchte RG]: unterteilt in feinblasig (eher
  Knistern)und grobblasig (Blubbern),
in ohrnah und ohrfern,
in klingend und nicht klingend
\end{description}
}{}{}{}{}

\Layout{2}{Atemexkursion}{%
Feststellung nicht seitengleicher Beatmung: Auflegen der
Hände und Beobachtung, ob sie sich beim Atmen seitengleich
verschieben
}{}{}{}{}

\Layout{2}{Stimmfremitus}{%
palpatorische Erfassung feiner Vibrationen mit den
Fingerspitzen über den Lungenpartien wenn der Patient mit
tiefer Stimme 99 spricht; bei vergleichender Palpation
verstärkte Vibration über Infiltrationen, vermindert über
Ergüssen und Pleuraschwarten
}{}{}{}{}

\Layout{2}{Bronchophonie}{%
Patient flüstert \Speech{\enquote{66}}: bei Abhören mit
Stethoskop ist das Geflüsterte besser hörbar über verdichtetem Lungengewebe,
Pneumonie, bei Infarkt, Atelektase, vermindert über
Pleuraerguss, -schwarte.
}{}{}{}{}










\subsection{Untersuchungsgang}




\Layout{2}{1.
Inspektion}{%

\begin{description}
  \item[Thoraxform] Beurteilung:
  Veränderungen der Thoraxform: Asymmetrien, Deformationen
  • Kyphose: Rundrücken (Wirbelsäulenkrümmung in sagittaler
  Ebene dorsal-konvex)
  • Lordose: Wirbelsäulenkrümmung ventral - konvex
  • Skoliose: Wirbelsäulenkrümmung in seitlicher Ebene
  leichte Brustkyphose und Lendenlordose: physiologisch!
  Kyphoskoliose: starke Verkrümmung der Wirbelsäule mit Atem- und
  Kreislaufbehinderung
  Veränderungen des Sternums:
  Trichterbrust /Pectus excavatus: unteres Sternum eingedellt, meist
  kongenital, Hühnerbrust (Pectus carinatus): Sternum ragt wie ein
  Kiel hervor, Keilbrust: rachitisch, kongenital, unbekannte Ursache
  bei Veränderungen des Sternums meist keine Atem-,
  Kreislaufbeeinträchtigung
  Fassthorax (Breite zu Tiefe 1:1):
  bei langdauernder obstruktiver Lungenerkrankung:
  Ausbildung einer überblähten Lunge (Emphysem), dabei sind die
  Supraclavicuklargruben oft vorgewölbt
  bei chronisch obstruktiver Lungenerkankung (COPD):
  „Blue bloater“: übergewichtig, cyanotischer Patient
  „Pink puffer“: eher magerer Patient mit rosa Haut
  \item[Atemfrequenz und Atemanstrengung]
  
  Atemfrequenz:
  normale Frequenz: 12-20 /min
  Tachypnoe >20 Atemzuge/min
  Bradypnoe: < 12/min
  Bei starker Atemnot versucht der/die Patient/In mit Hilfe der
  Atemhilfsmuskulatur im Sitzen mit aufgestützten Armen ein größeres
  Atemvolumen zu erreichen (auxiliäre Muskulatur).
  Pathologische Atmungstypen:
  • Kussmaulsche Atmung: vertiefte Atmung unregelmäßiger
  Frequenz; bei metabolischer Azidose, diabetischem Coma
  • • Cheyne-Stokes Atmung: periodisches An- und Abschwellen der
  Atemtiefe; Hirnschädigung, Medikamente (Drogen!)
  • • Biot'sche Atmung: unregelmäßige, terminale
  Schnappatmung; Hirntraumen, Hirndrucksteigerung
  Stridor:
  in- oder exspiratorisches auf Distanz hörbares Stenosegeräusch bei
  Verengung der großen
  Atemwege
  • inspiratorisch laryngealer Stridor: bei Entzündung (obstruktive
  Laryngobronchitis –
  Pseudokrupp) oder Fremdkörper (Kinder!)
  • trachealer Stridor: oft in- und exspiratorisch bei Verengung durch
  Struma oder Tumor
  • bronchialer exspiratorischer Stridor: bei Asthma
  \item[Hautfarbe]
\end{description}

}{}{}{}{}








\Layout{2}{2. Palpation}{%

\begin{description}
  \item[Atemexkursion]
  \item[Stimmfremitus]
\end{description}
}{}{}{}{}





\Layout{2}{3. Perkussion}{%

\begin{description}
  \item[vergleichende Perkussion] der
  Lungenfelder
  
  Klopfschall kann durch Adipositas, Emphysem, Thoraxasymmetrien
  verfälscht sein
  \item[abgrenzende Perkussion] der
  Lungengrenzen
  
  basal: Lungen/Lebergrenze mit leiser Perkussion in der
  Mediaklavikularlinie von sonorem
  Lungenschall abwärts perkutieren bis Schenkelschall etwa am
  Oberrand der 6. Rippe auftritt; bei stärkerer Perkussion wird die
  Zwerchfellkuppe erfasst und damit die Grenze zu hoch
  angesetzt; Grenzen: in der mittleren Axillarlinie 8. Rippe,
  Skapularlinie 9. Rippe, paravertebral Dornfortsatz des 11 BWK
  \item[Atemverschieblichkeit]
  (Erst in Exspiration, dann in Inspiration von oben nach unten
  perkutieren) normal etwa 4-6 cm; Lungenbasis bei Inspiration in
  Skapularlinie bis 5 QF (Querfinger), in
  Paravertebralinie 3 QF verschieblich;
  verringert z.\,B. bei Pleuraverwachsungen nach Pleuritis, Erguss;
  tiefstehend mit verringerter Verschieblichkeit bei Emphysem und
  obstruktiven und restriktiven Ventilationsstörungen
  \item[Krönigschen Schallfelder]Breite
  bestimmen
  
  In Schulterhöhe die Krönig`schen Schallfelder über Lungenspitzen
  etwa 3 QF breit, eingegrenzt bzw. nicht auffindbar bei Prozessen an
  der Lungenspitze (Tb), bzw. breiter bei Emphysem
\end{description}
}{}{}{}{}


\Layout{2}{4.
Auskultation}{%

\begin{description}
  \item[Lungenfelder im Seitenvergleich] (Patient atmet tief
  mit offenem Mund)
  
  Vesikuläratmen: normales Atemgeräusch, über Lungengewebe etwas
  leiser in der Exspiration
  • Bronchialatmen: wie über Trachea oder oberen Lungenpartien
  mit großen Bronchien, z.\,B. auch über rechtem Apex, verschärft
  in In- und Exspiration über verdichtetem Lungengewebe z.\,B.
  Entzündung (Pneumonie)
  • vermindertes Atemgeräusch: einseitig über Ergüssen,
  Pleuraschwarte und fast ganz bei Atelektasen, beidseitig bei
  Emphysem
  • asthmatisches Atmen: Exspirium betont und verlängert, meist
  mit Giemen und Brummen
  • trockene Rasselgeräusche (RG): Giemen, Brummen, Pfeifen;
  durch zähen Schleim; z.\,B. Bronchitis – chronische
  Raucherbronchitis
  • feuchte Rasselgeräusche: Zeichen von Flüssigkeitsansammlung
  in Bronchien und Alveolen; fein- bis grobblasige Rasselgeräusche
  (RG), Knisterrasseln: bei Pneumonie, Lungenstauung – ohrnah
  (z.\,B. Pneumonie), ohrfern (z.\,B. Herzinsuffizienz)
  \item[Bronchophonie] testen
\end{description}

}{}{}{}{}
5.
Dokumentation







































\subsection{Typische Befunde}
Beispiele: I = Inspektion, P = Palpation, Pk = Perkussion, A = Auskultation,

\Layout{2}{Bronchitis}{%
\begin{labeling}[:]{mm}
  \item[I] trockener Husten
%   \item[P] 
%   \item[Pk] 
%   \item[A] 
%   \item[»] 
\end{labeling}
}{}{}{}{}

\Layout{2}{Asthma bronchiale (Anfall)}{%
\begin{labeling}[:]{mm}
  \item[I] Dyspnoe, beschleunigte Atmung, sichtbare
  Behinderung der Exspiration, Atemhilfsmuskulatur
%   \item[P] 
  \item[Pk] (nur im Anfal) hypersonor, tiefstehende Basen
  \item[A] verlängertes Exspirium, trockene, giemende, pfeifende Rasselgeräusche
%   \item[»] 
\end{labeling}
}{}{}{}{}

\Layout{2}{Chronische Bronchitis (Raucher)}{%
\begin{labeling}[:]{mm}
  \item[I] Verlauf mit intermittierenden Verschlechterungen – Dyspnoe, Husten
%   \item[P] 
  \item[Pk] zunehmend hypersonorer Klopfschal, zunehmende Lungenblähung
  \item[A] trockene RG, verlängertes Exspirium
%   \item[»] 
\end{labeling}
}{}{}{}{}

\Layout{2}{Emphysem}{%
\begin{labeling}[:]{mm}
  \item[I] Fassthorax, horizontaler Rippenverlauf
%   \item[P] 
  \item[Pk] hypersonorer KS, Basen tiefstehend, unverschieblich
  \item[A] abgeschwächtes Atemgeräusch, verlängertes Exspirium, trockene Rasselgeräusche
%   \item[»] 
\end{labeling}
}{}{}{}{}

\Layout{2}{Pneumonie (bakteriell)}{%
alles auch abhängig von Größe und Lage
\begin{labeling}[:]{mm}
  \item[I] Fieber, evtl. beschleunigte Atmung, Husten, bei schweren Fällen Dyspnoe
  \item[P] Stimmfremitus (99) verstärkt , evtl. abgeschwächte Atemexkursionen
  \item[Pk] evtl. Dämpfung abhängig von Größe, Abstand zur Oberfläche (<5cm)
  \item[A] Bronchialatmen, Knisterrasseln, feinblasige
  Rasselgeräusche, Bronchophonie (66)
%   \item[»] 
\end{labeling}
}{}{}{}{}
\Layout{2}{Atypische Pneumonie}{%
(Viren, Mycoplasmen, Pilze, Legionela)
\begin{labeling}[:]{mm}
%   \item[I] 
%   \item[P] 
%   \item[Pk] 
%   \item[A] 
  \item[»] nur geringe Lokalbefunde, massivere Veränderungen im Röntgen
\end{labeling}
}{}{}{}{}

\Layout{2}{Atelektase}{%
\begin{labeling}[:]{mm}
  \item[I] asymmetrische Atmung mit inspiratorischer Einziehung der Zwischenrippenräume, nur bei
%   \item[P] 
%   \item[Pk] 
%   \item[A] 
%   \item[»] 
\end{labeling}
}{}{}{}{}

\Layout{2}{Verschluss von Bronchien 2. oder 3. Ordnung}{%
\begin{labeling}[:]{mm}
  \item[I] 
  \item[P] Stimmfremitus abgeschwächt
  \item[Pk] evtl. leichte Dämpfung, abhängig vom Ausmaß
  \item[A] Atemgeräusch und Bronchophonie abgeschwächt
  \item[»] 
\end{labeling}
}{}{}{}{}

\Layout{2}{Bronchialkarzinom}{%
\begin{labeling}[:]{mm}
%   \item[I] 
%   \item[P] 
%   \item[Pk] 
%   \item[A] 
  \item[»] lange Zeit ohne Veränderung des physikalischen Befundes
\end{labeling}
}{}{}{}{}

\Layout{2}{Pneumothorax}{%
\begin{labeling}[:]{mm}
  \item[I] Interkostalräume vorgewölbt
  \item[P] Atemexkursion abgeschwächt, Stimmfremitus abgeschwächt
  \item[Pk] hypersonorer bis tympanitischer Klopfschal, evtl. Schachtelton
  \item[A] abgeschwächte bis fehlende Atemgeräusche und Bronchophonie
  \item[»] (Cave: Spannungspneumothorax, absoluter Notfall:
  Punktion 4./5. ICR mittlere Axillarlinie)
\end{labeling}
}{}{}{}{}

\Layout{2}{Pleuritis sicca}{%
\begin{labeling}[:]{mm}
  \item[I] evtl. Schonhaltung, geringere Atemexkursionen, Patient liegt auf kranker Seite
  \item[P] Atemexkursionen abgeschwächt, Nachhängen der kranken Seite
  \item[Pk] verminderte Basenverschieblichkeit
  \item[A] trockene Reibegeräusche, „Lederknarren“, Verschwinden bei Auftreten eines Ergusses
%   \item[»] 
\end{labeling}
}{}{}{}{}

\Layout{2}{Pleuraerguss}{%
\begin{labeling}[:]{mm}
  \item[I] bei großem Erguss evtl. Nachschleppen der kranken Seite
  \item[P] Atemexkursionen abgeschwächt, evtl. asymmetrisch, Stimmfremitus abgeschwächt
  \item[Pk] Dämpfung, Ellis-Damoisau'sche Linie,
Atemverschieblichkeit aufgehoben, im Sitzen nach oben abnehmende Dämpfung
  \item[A] Atemgeräusch abgeschwächt
%   \item[»] 
\end{labeling}
}{}{}{}{}

\Layout{2}{Pleuraschwarte}{%
\begin{labeling}[:]{mm}
  \item[I] evtl. Einziehung des Interkostalraumes
  \item[P] Stimmfremitus (99) vermindert, abgeschwächte Atemexkursion
  \item[Pk] KS gedämpft
  \item[A] abgeschwächte Atemgeräusche
%   \item[»] 
\end{labeling}
}{}{}{}{}

\Layout{2}{Lungenstauung bei kardialer Links-Dekompensation}{%
\begin{labeling}[:]{mm}
  \item[I] Dyspnoe, Orthopnoe (Luftnot beim flachen Liegen, die sich durch aufrechte Körperhaltung
bessert)
%   \item[P] 
  \item[Pk] eher gedämpft
  \item[A] feinblasige Rasselgeräusche über Basis, bis zunehmend über ganzer Lunge
%   \item[»] 
\end{labeling}
}{}{}{}{}



\section{Herz}

Vorwissen
Für folgende Begriffe schreiben Sie bitte eine kurze Definition in Ihr Skript:

Eupnoe
Dyspnoe
Orthopnoe
Zyanose zentrale/periphere
Facies mitralis
Periphere Ödeme
Lymphödem
Jugularvenenpuls
Herzspitzenstoß (HSS)
Erb'scher Punkt wo
Musset'sches Zeichen
Auskultationspunkt der Aortenklappe
Auskultationspunkt der Mitralklappe
Auskultationspunkt der Pulmonalklappe
Auskultationspunkt der Trikuspidalklappe
Tachykardie
Bradykardie
Arrhythmia absoluta
respiratorische Arrhythmie
Spaltung des 2. HT
Mitralklappenöffnungston
Ejection click
Funktionelles Herzgeräusch
Pericarditis
Wasserhammerpuls
Aortenisthmusstenose
Subclavian-Steal-Syndrome
AtriumSeptumDefekt
VentrikelSeptumDefekt
Persistierender Duktus Botalli






Die Herzaktion lässt sich in 4 Phasen teilen:
Anspannungs-, Austreibungs-
\
/
Systole
Entspannungs-, und Füllungsphase.
\
/
Diastole

Der 1. Herzton (HT) ist niedrigfrequent und entsteht in der
Anspannungsphase durch Anspannung der Ventrikelmuskulatur,
Schluss der AV-Klappen und Schwingung des ganzen Systems. Er ist
carotis-pulssynchron hörbar.

Der 2. Herzton ist etwas höher und entsteht in der
Entspannungsphase durch Schluss der Taschenklappen und
Schwingung der Systems.
Auskultatorisch ist die Systole also die Phase zwischen 1. und 2. HT.
Jetzt ist der Puls tastbar!
Die Diastole damit die Phase zwischen 2. und 1. HT.

Herzgeräusche werden so in Systolika und Diastolika unterschieden.
Ihre Lautstärke wird in Sechstel angegeben:

1: leise, kaum hörbar
2: leise, gut hörbar
3: mittelaut
4: laut
5: sehr laut (fortgeleitet auch außerhalb der Herzgegend zu hören)
6: Distanzgeräusch, mit freiem Ohr zu hören

Weitere Charakterisierung:
zeitliche Zuordnung: systolisch: prä-, spät-, holodiastolisch:
früh-, meso-, protodiastolisch (meist pathologisch)
zeitlicher Verlauf: crescendo
decrescendo
gleichförmig, bandförmig, spindelförmig
Charakter: Frequenz: hoch-, niederfrequent
Tonhöhe: z.\,B. hell, laut
Ausstrahlung: Axilla, Karotiden
Veränderung durch Atmung, Lage:
z.\,B. physiologische Spaltung des 2. Herztones, Umlagerung
oder Bewegung, z.\,B. akzidentelle Geräusche

Auskultation:
Auskultationsstellen der Herzklappen:
• Erb'scher Punkt: 3/4. ICR parasternal links (grundsätzlich
alle Töne und Geräuschehier gut hörbar)
• Aortenklappe: 2. ICR rechter Sternalrand
• Mitralklappe, Herzspitze: 5. ICR medial der MCL
• Trikuspidalklappe: 4/5. ICR rechts parasternal
• Pulmonalklappe: 2. ICR links parasternal

Rhythmus:
regelmäßiger Herzschlag mit Frequenz um etwa 60-80 /min
Die bei jungen (und bei „vegetativ“ übererregbaren) Menschen gefundene respiratorische
Arrhythmie mit Frequenzzunahme bei Inspiration und Verringerung bei Exspiration hat
keinen Krankheitswert.











Untersuchungsgang
1.
Freimachen
des Patienten
Unter Wahrung der Intimsphäre
Bestimmung der Atemfrequenz
2.
Inspektion

Begutachtung der Thoraxform
Erkennen einer Zyanose
Bestimmung des Jugularvenenpuls
Achten auf periphere Ödeme

3.
Palpation
Bestimmung des Herzspitzenstoßes(HSS)


4.
Auskultation
bei gleichzeitiger Palpation des Carotes-
Pulses wird der 1. Herzton bestimmt und
ein etwaiges Herzgeräusch der Systole
bzw Diastole zugeordnet und zu seinem
Ort maximaler Lautstärke (punctum
maximum=p.m.) verfolgt. Die meisten
Herzgeräusche sind in maximaler
Expiration am Besten zu hören, da das
Herz sich der Brustwand nähert
auf ein peripheres Pulsdefizit achten
\begin{itemize}
  \item Auskultation:
  \begin{itemize}
    \item Erb'schen Punkts
    \item p.m. der Aortenklappe
    \item p.m. der Pulmonalklappe
    \item p.m. der Trikuspidalklappe
    \item p.m. der Mitralklappe
    \item linke Axilla
    \item beide Carotiden (nacheinander)
  \end{itemize}
  \item Beurteilen von:
  \begin{itemize}
    \item Rhythmisch?
    \item Herzfrequenz?
    \item Qualität der Herztöne ?
    \item Herzgeräusche beschrieben bezügl. Zeit,
    Lautstärke, p.m., Fortleitung und
    Charakter
  \end{itemize}
\end{itemize}






5.
Dokumentation






Zusatzinformation, mögliche Befunde und deren Interpretation

normal, Dyspnoe: Belastungsdyspnoe, Ruhedyspnoe, Orthopnoe
Thoraxform verändert: starke Wirbelsäulenveränderungen können zu
einer Rechtsherzbelastung führen; siehe Thorax-Inspektion
zentrale/periphere Zyanose (Facies mitralis)
Halsvenenstauung (obere Einflussstauung, Jugularvenenpuls)
Ödeme an Unterschenkeln, Knöcheln

Herzspitzenstoß (HSS): normal im 5. ICR innerhalb der MKL, 2 QF breit
HSS verlagert:
• nach außen: bei Linksherzvergrößerung, evtl. verbreitert und
hebend
• nach oben: bei Rechtsherzvergrößerung
• verbreitert: Hypertrophie
• nicht tastbar: bei Fettleibigkeit, Zwerchfellhochstand, Emphysem
• schwach: z.\,B. bei Mitralstenose
• schwirrend: bei Vorhof-,Ventrikelseptumdefekt
•
•

organisch: Klappendestruktion, Verkalkung, Hypertrophie
funktionell: bei Änderung des Blutvolumens oder des Herz-
Minutenvolumens, immer systolisch, oft bei Jugendlichen, keiner
Struktur zuordenbar, nie diastolisch, kein Krankheitswert
regelmäßiger Herzschlag mit Frequenz um etwa
60-80 /min
Die bei jungen (und bei „vegetativ“ übererregbaren) Menschen
gefundene respiratorische Arrhythmie mit Frequenzzunahme bei
Inspiration und Verringerung bei Exspiration hat keinen Krankheitswert.



\subsubsection{Typische Befunde und Interpretation}

\Layout{2}{1. Herzton}{%
\begin{labeling}[:]{mm}
%   \item[I] 
%   \item[P] 
%   \item[Pk] 
  \item[A] laut: bei „Stress“ (hyperdyname Kreislaufsituation), Fieber, Anämie, Gravidität,
Hyperthyreose

paukend: bei Mitralstenose

leise: bei Herzinsuffizienz, Mitralinsuffizienz, evtl. bei Vorhofflimmern, Myocarditis
gespalten: bei Links- und Rechtschenkelblock, Extrasystolen
  \item[»] 
\end{labeling}

}{}{}{}{}


\Layout{2}{2. Herzton}{%
\begin{labeling}[:]{mm}
%   \item[I] 
%   \item[P] 
%   \item[Pk] 
  \item[A] laut: bei Hypertonie, Aortensklerose
leise bis fehlend: Aortenstenose

physiologische Spaltung: bei Inspiration verstärkt - Aortenklappe schließt vor
Pulmonalklappe - bei Jugendlichen in Inspiration > 0,03 sek., verschwindet in Exspiration

paradox gespalten: Pulmonalklappe schließt vor Aortenklappe: bei Aortenstenose,
Aorteninsuffizienz, Linksschenkelblock

fixierte Spaltung: möglich bei Volumsbelastung durch links-rechts Shunt (z.\,B.
Vorhofseptumdefekt), pulmonaler Hypertonie
  \item[»] 
\end{labeling}

}{}{}{}{}



\Layout{2}{3. Herzton}{%:
\begin{labeling}[:]{mm}
%   \item[I] 
%   \item[P] 
%   \item[Pk] 
  \item[A] pm: Herzspitze, ventrikulärer Füllungston in früher Diastole, am besten mit Schalltrichter zu
hören; 3. HT und leiser 1. und 2. HT klingt wie Pferdegalopp (früher protodiastolischer
Galopp genannt)

  \item[»] Vorkommen: bei rascher Füllung oder verminderter Dehnbarkeit bei Mitralinsuffizienz,
Herzinsuffizienz, Kardiomyopathie, Myocardinfarkt
bei Kindern und Jugendlichen physiologisch bei hyperdynamer Kreislaufsituation - aber auch
bei Tachycardie, Fieber, Sportlern, und in der Schwangerschaft
\end{labeling}

}{}{}{}{}

\Layout{2}{4. Herzton}{%:
\begin{labeling}[:]{mm}
%   \item[I] 
%   \item[P] 
%   \item[Pk] 
  \item[A] kräftige Vorhof-Kontraktion (Vorhofgalopp, präsystolischer Galopp), pm Erbscher Punkt;
knapp vor dem 1. HT
  \item[»]  bei hypertropher Kammermuskulatur (Linksherzhypertrophie),
Hypertonus, Aortenstenose, Myocardinfarkt
bei Jugendlichen physiologisch direkt vor Systole; bei verstärkte Kontraktion eines
hypertrophen Vorhofes, und/oder Aufprall des ausgestoßenen Vorhofblutes auf eine
hypertrophierte Ventrikelwand, bei Linkshypertrophie, bei arterieler Hypertonie oder
Aortenstenose
\end{labeling}

}{}{}{}{}

\Layout{2}{Ejection click}{%Weitere Herztöne:
\begin{labeling}[:]{mm}
%   \item[I] 
%   \item[P] 
%   \item[Pk] 
  \item[A] Dehnungston bei poststenotischer Erweiterung von Aorta und Pulmonalis;
PM 2. ICR rechts/links parasternal
Mitraler Öffnungston bei Mitralstenose - MÖT
frühsystolischer Austreibungston bei Aorten- oder Pulmonalstenose
  \item[»] 
\end{labeling}

}{}{}{}{}

\Layout{2}{Mitralklappenstenose (MS)}{%:
\begin{labeling}[:]{mm}
  \item[I] event. Facies mitralis
  \item[P] Herzspitzenstoß (HSS) abgeschwächt – nicht fühlbar, (linker Ventrikel füllt sich schlecht)
  \item[Pk] 
  \item[A] 1. HT: paukend (verzögerte Kammerfüllung, ruckartiger Klappenschluss – evtl. in
entgegengesetzte Richtung – durch die Kammerkontraktion), Systole: frei, 2. HT: normal,
knapp gefolgt vom mitralen Öfnungston (MÖT) – je kürzer der Abstand, desto höhergradiger
die Stenose – dem ein diastolisches, raues, niederfreqentes, leises Decrescendo folgt; wenn
noch Sinusrhythmus präsystolisches Crescendo
  \item[»] 
\end{labeling}
}{}{}{}{}

\Layout{2}{Mitralklappeninsuffizienz (MI)}{%
\begin{labeling}[:]{mm}
%   \item[I] 
  \item[P] HSS verbreitert, nach links und unten verlagert, hebend
%   \item[Pk] 
  \item[A] 1. HT: leise (Klappenschluss der Mitralklappe abgeschwächt); dann holosystolisches,
bandförmiges, gießendes Geräusch, PM Herzspitze fortgeleitet in Axilla; 2. HT: gespalten; 3.
HT: hörbar bei erheblicher Volumenüberlastung
  \item[»] 
\end{labeling}
}{}{}{}{}

\Layout{2}{Aortenklappenstenose (AS)}{%:
\begin{labeling}[:]{mm}
  \item[A] 1. HT: normal bis abgeschwächt, lautes, raues,
  spindelförmiges systolisches Austreibungsgeräusch, PM 2. ICR rechts fortgeleitet in die Carotiden; 2. HT: abgeschwächt und
verspätet, evtl. paradox gespalten; eventuell 4. HT: kräftige Vorhofkontraktionen; 3. HT: hörbar
bei kardialer Dekompensation
%   \item[I]
  \item[P] HSS hebend; Pulsus parvus et tardus (kleiner
  Puls, träger Anstieg der Pulswelle), systolisches Schwirren links bei Vornüberbeugen
%   \item[Pk]
%   \item[»]
\end{labeling}
}{}{}{}{}

\Layout{2}{Aortenklappeninsuffizienz (AI)}{%:
\begin{labeling}[:]{mm}
  \item[A] 1. HT: unaufällig, leise, (evtl. ejection click),
  2. HT: abgeschwächt, leises, gießendes diastolisches Decrescendo, PM 2. ICR rechts, Erb'scher Punkt
  \item[I] pulsierende Gefäße bei großem Schlagvolumen
  (A.carotis)
  \item[P] HSS nach links und unten verlagert, verbreitert
  und hebend, Pulsus celer et altus – Wasserhammerpuls
%   \item[Pk]
%   \item[»]
\end{labeling}
}{}{}{}{}


\Layout{2}{Pulmonalstenose}{%:
\begin{labeling}[:]{mm}
  \item[A] syst. Austreibungsgeräusch, fixierte Spaltung 2. HT, PM Pulmonalregion (2. ICR links),
Ejection click
%   \item[I] 
%   \item[P] 
%   \item[Pk] 
%   \item[»] 
\end{labeling}
}{}{}{}{}


\Layout{2}{Trikuspidalinsuffizienz}{%
\begin{labeling}[:]{mm}
  \item[A] holosystolisches, bandförmiges Geräusch, PM 4. ICR rechts parasternal
  \item[I] Pulsation der Halsvenen (positiver Venenpuls)
  \item[P] Halsvenenpuls
  \item[Pk] Herzdämpfung nach rechts
%   \item[»] 
\end{labeling}
}{}{}{}{}



\Layout{2}{Aortenisthmusstenose}{%:
\begin{labeling}[:]{mm}
  \item[A] systolisch-diastolisches Geräusch mit PM ICR
  parasternal links und im Rücken
%   \item[I]
  \item[P] Pulse der oberen Extremität kräftig, der unteren
  schwach, RR Diferenz OE > UE
%   \item[Pk]
%   \item[»]
\end{labeling}
}{}{}{}{}

\Layout{2}{Ventrikelseptumdefekt}{%:
\begin{labeling}[:]{mm}
  \item[A] holosystolisches, spindelförmiges Geräusch, laute
  Herztöne, 2. HT gespalten mit PM 3. ICR links
%   \item[I]
%   \item[P]
%   \item[Pk]
%   \item[»]
\end{labeling}
• }{}{}{}{}

\Layout{2}{Vorhofseptumdefekt}{%:
\begin{labeling}[:]{mm}
  \item[A] meist Systolikum, eventuell Diastolikum, fixierte
  Spaltung des 2. HT
%   \item[I]
  \item[P] Schwirren präcordial
%   \item[Pk]
%   \item[»]
\end{labeling}
}{}{}{}{}

\Layout{2}{Perikardgeräusche}{%:
\begin{labeling}[:]{mm}
%   \item[I] 
%   \item[P] 
%   \item[Pk] 
  \item[A] herzsynchrones, ohrnahes Reibegeräusch bei trockener Pericarditis, auch Systole/Diastole
überlagernd, oft flüchtig, verschwindet bei Auftreten eines Ergusses
%   \item[»] 
\end{labeling}

}{}{}{}{}










\section{Abdomen und rektale Untersuchung}
Vorwissen

Akutes Abdomen
Xiphoid
direkte Inguinalhernie
indirekte Inguinalhernie
Ascites
Meteorismus
Rektusdiastase
Striae
paralytischer Ileus
obstruktiver Ileus
Peristaltik
Abwehrspannung
Appendizitis
McBurney Punkt
Lanz'scher Punkt
Blumberg Zeichen
Psoaszeichen
Obturatorzeichen
Divertikulitis
Zystennieren
Cushing Syndrom
Leberzirrhose
Fettleber
Caput medusae
Courvasier'sches Zeichen
Nierenarterienstenose
Flankendämpfung
Undulationsphänomen
Marisquen
Hämorrhoiden
Meläna
Prostatitis
Kohlrau'sche Falte

\subsubsection{Kurzer theoretischer Überblick}


Das Abdomen wird in 4 Quadranten eingeteilt, die gedachten
Trennlinien gehen hier jeweils durch den Nabel. Die
Entspannung des Patienten ist bei der Untersuchung des
Abdomens besonders wichtig, deshalb wird der Patient mit
einem Kissen unter dem Kopf auf dem Rücken gelagert,
eventuell können die Beine angewinkelt werden. Bei der
Untersuchung mit dem Patienten reden, ablenken, unterhalten,
tief ein- und ausatmen lassen, Untersuchung mit warmen
Händen, kalte Hände führen zu einer Abwehrspannung; Abdomen
systematisch abtasten. Vor der Palpation wird eine
ausführliche Auskultation in allen 4 Quadranten vorgenommen.
Die Palpation wird erst oberflächlich und dann tief
(beidhändig) vorgenommen, hierbei bewegt man sich vorsichtig
auf den Punkt des möglichen Schmerzes zu. Abwehrspannung
kann man am besten beim Einatmen beurteilen.

Zu erfassende Organe: Magen, Leber, Gallenblase, Milz,
Nieren, Pankreas, Dünn- und Dickdarm, Gefäße, Harnblase und
Genitalorgane

Ascites ist eine Flüssigkeitsansammlung in der Bauchhöhle.
In Rückenlage perkuttiert man ein Flankendämpfung - beim
Drehen des Patienten in rechte und linke Seitenlage eine
wechselnde Dämpfung. Ausserdem kann bei ausgedehntem Ascites
ein Undulationsphänomen auftreten: eine Hand an einer Seite
am Abdomen anlegen, auf der gegenüberliegenden Seite mit den
Fingern anklopfen - bei Vorliegen einer Flüssigkeitsfüllung
verspürt man das Ankommen eines Stoßes

Auskultation:

Kratzauskultation (unverlässlich): Aufsetzen des Stethoskops
im Epigastrium über der vergrößerten Leber - mit der
Fingerspitze über Leberareal streichen – das dabei hörbare
Geräusch verschwindet bei Verlassen des Leberbereichs






\subsubsection{Untersuchungsgang}



\Layout{2}{1.Inspektion}{%

Verhältnis des Abdomens zum Thoraxniveau
Abdomen über dem Thoraxniveau:
• bei Adipositas, Meteorismus, Aszites, großen Tumoren, überfüllter
riesiger Harnblase
• bei Harnverhalten, Zystennieren und Gravidität
• Eingezogenes Abdomen:
• bei Kachexie, evtl. Meningitis
• Rektusdiastase:
• beim Aufrichten aus liegender Position Auseinanderklaffen der
beiden Hälften des Musc. recti bei Multipara und bei Aszites


Achten auf Haut

Hautbeschaffenheit: trocken, ödematös, adipös, kachektisch, Exsikkose
Behaarungsabnormitäten: weiblicher Behaarungstyp (Bauchglatze), z.\,B.
bei Leberzirrhose

Striae: z.\,B. rötlich bei Cushing Syndrom, weißlich nach Schwangerschaft
Kollateralkreislauf (Caput medusae)
<<<BILD>>>

Achten auf Leberhautzeichen

Achten auf Rektusdiastase

Achten auf Narben

Achten auf Hernien

Leistenhernie (direkt/indirekt), Femoralhernie, Narbenhernie,
epigastrische Hernie (unter Schwertfortsatz)
}{}{}{}{}

\Layout{2}{2.
Auskultation}{%


Langes Abhören des Abdomens in allen 4
Quadranten (achten auf Darm- u.
Gefäßströmungsgeräusche)

Peristaltik:
• normal: eher leise, diskontinuierlich
• gesteigert: bei Durchfall, Magen/Darmblutung, Obstruktion
• fehlend: bei paralytischem Ileus (Totenstille)
Strömungsgeräusche:
von höhergradig stenosierten Arterien: Nierenarterienstenose, Aa.
iliacae, Aa. femorales

}{}{}{}{}


\Layout{2}{3.
Palpation}{%


Des Leistenkanals

Der inguinalen LK

Des Abdomens erst oberflächlich dann tief
(achten auf auf Druckschmerz,
Abwehrspannung, Resistenzen,
Raumforderungen und die Patientenreaktion)

Als Normalbefund sind kein abdominaler Druckschmerz und keine
Resistenzen festzustellen
Schmerzpunkte: bei lokalen Entzündungen, Auslassschmerz, z.\,B.
McBurney bei Appendizitis
Schmerzen:

bei lokalisiertem Druckschmerz – Quadranten oder Region angeben

• rechter oberer Quadrant: z.\,B. Cholezystitis, Pankreatitis,
Ulkusperforation

• rechter unterer Quadrant: z.\,B. Appendizitis -
McBurney'scher und Lanz'scher Punkt und Nabel

• linker oberer Quadrant: z.\,B. Pankreatitis, Milzinfarkt, Nierenstein,
Ureterstein

• linker unterer Quadrant: z.\,B. Divertikulitis, Nieren-, Ureterstein,
Adnexitis (bds. möglich), Loslasschmerz – lokale Peritonitis

Abwehrspannung:
geringer Druck führt durch Schmerzauslösung zur Spannung der
Bauchdecke (Zeichen für lokale Peritonitis), bretthartes (akutes)
Abdomen

Untersuchung auf Ascites mit
Stosswellenperkussion
Zusatzinformation, mögliche Befunde und deren Interpretation

Testen auf Ascites:
DD: Adipositas, Blähungen, Stuhlmassen

Inspektion:
bei größeren Ascitesmengen: Flanken wölben sich vor, massiver Ascites
halbkugelig vorgewölbt, Nabel verstrichen, evtl. Nabelhernie, dünne
Bauchdecke; Venenzeichnung an der lateralen Bauchwand (Caput
medusae): bei Patienten mit portaler Hypertension (Kollateralen
zwischen Pfortader und V. cava)



Untersuchung auf Flankendämpfung
(Perkussion)
}{}{}{}{}














\Layout{2}{4.
Leber und Gallenblase}{%



Leberränder mittels Kratzauskultation und
Perkussion bestimmen (die Palpation erfolgt
in Inspiration in der Medioklavikularlinie, falls
dort nicht tastbar, medial und lateral davon)

Palpation: bei Inspiration tritt Leber tiefer, in rechter MCL kann der
Leberrand 1-2 QF unter
dem Rippenbogen palpierbar sein:
weich, glatt, keine Eigenpulsation, nicht druckempfindlich, der
Leberrand ist scharfrandig, glatte Oberfläche;
Leberzirrhose oder evtl. Metastasen: plumper Leberrand, erhöhte
Konsistenz, gebuckelte Oberfläche;
bei Aszites Verwendung der Stoßpalpation (aufgestellte Finger unter
Ribo)

Perkussion: Bestimmung der Lungen-Lebergrenze, Untergrenze im
Abdomen schwer fassbar;
}{}{}{}{}


\Layout{2}{5. Milz}{%

Bimanuelle Palpation in tiefer Inspiration
(normale Milz ist nicht tastbar)
}{}{}{}{}


\Layout{2}{6. Nieren}{%

Bimanuelle Palpation in tiefer Inspiration
(eventuell unterer Pol tastbar)

Palpation: Nieren normalerweise nicht palpabel; um eine Vergrößerung
zu erfassen palpiert man bimanuell, wobei man mit einer Hand vom
Rücken nach ventral drückt, mit der anderen Hand unter dem
Rippenbogen bei Inspiration entgegendrückt; bei Normalbefund evtl.
rechter unterer Nierenpol erreichbar (linke Niere steht höher);
Vergrößerung bei Nierenzysten, Tumoren;
durch Beklopfen des Rückens in der Nierengegend: Feststellung, ob
Nierenlager klopfdolent (u.\,U. bei entzündlichen Prozessen –
Pyelonephritis, Abszess)

Überprüfung der Klopfempfindlichkeit der
Nierenlager
}{}{}{}{}

\Layout{2}{7.
Harnblase}{%


Bestimmung der Größe und Füllung durch
Perkussion
}{}{}{}{}


\Layout{2}{8.
Rektale
Untersuchung}{%



palpatorische Beurteilung der Analregion
palpatorische Beurteilung des Rektums
palpatorische Beurteilung der Prostata


Die rektale Untersuchung: sollte bei jeder eingehenden klin. phys.
Untersuchung durchgeführt
werden. Sie kann Hinweise auf Tumore (\VonBis{20}{30}\PC*
aller Dickdarmkarzinome sind mit dem Finger erreichbar und palpabel) und
Erkrankungen von Prostata und Uterus geben;
mögliche pathologische Befunde:

• reduziertem Sphinkertonus: altersbedingt, Analprolaps, Proktitis,
neurogene Störung
• peritoneale Reizung: z.\,B. Peritonitis, Adnexitis, Appendizitis
• Rektumkarzinom
• Hämorhoiden
• Polypen
• Obstipation (harte Kotballen)
• Myomen des Uterus
• Entzündung im Bereich der Parametrien
• Prostatahypertrophie, Prostatakarzinom, Prostatitis
}{}{}{}{}
